\documentclass[10pt,twocolumn]{article}
\usepackage{graphicx}
\usepackage{amsmath}
\usepackage{amssymb}
\usepackage{cite}
\usepackage{array}
\usepackage{booktabs}
\usepackage{multirow}
\usepackage{hyperref}
\usepackage[margin=0.75in]{geometry}
\usepackage{xcolor}
\usepackage{float}

% IEEE style formatting
\renewcommand{\baselinestretch}{1.0}
\pagestyle{plain}

\title{Unified Edge-to-Cloud Evaluation Framework: Benchmarking IoT Data Reduction Strategies on Synthetic and Real-World Datasets}

\author{Aman Sharma$^{*}$, Dr. Madhurima Buragohain\\
Department of Computer Science and Engineering\\
Dr. B. R. Ambedkar National Institute of Technology, Jalandhar, Punjab\\
\textit{*Corresponding author: 24203003@nit-j.ac.in}}

\date{}

\begin{document}

\maketitle

\begin{abstract}
Internet of Things (IoT) systems generate massive volumes of time-series data that must be efficiently transmitted from edge devices to cloud servers. While numerous data reduction, communication protocol, and serialization strategies exist, prior evaluations examine these techniques in isolation, preventing holistic understanding of efficiency-accuracy trade-offs. Furthermore, most evaluations use only synthetic data, lacking validation on real-world sensor deployments. This paper presents a unified edge-to-cloud evaluation framework that jointly benchmarks four data reduction strategies (Adaptive Sampling, Adaptive Threshold, Aggregation, Filtering), two protocols (HTTP, MQTT), and two serialization formats (JSON, CBOR) under identical experimental conditions on \textit{both synthetic and real-world datasets}. Through 648 systematic experiments on synthetic data with 12 runs per configuration, followed by validation on the Intel Berkeley Research Lab dataset (54 real sensors, 2.3M+ readings), we demonstrate comprehensive performance characterization. On synthetic data, we achieve up to 99.14\% data reduction using chained strategies while maintaining RMSE $\leq$ 0.105. On real-world Intel Lab data, AdaptiveThreshold achieves 52.64\% reduction (vs. 93.11\% on synthetic), demonstrating realistic production-grade performance expectations. Our findings show that CBOR serialization reduces message size by 54\% compared to JSON without impacting accuracy, and that protocol choice (HTTP vs. MQTT) does not affect reduction metrics. Real-world validation reveals that sensor drift, missing values, and unpredictable patterns reduce achievable reduction by approximately 40\% compared to synthetic data, a critical insight for IoT practitioners. This work provides data-driven guidelines for selecting reduction-protocol-serialization combinations tailored to specific constraints, validated against real-world sensor data.
\end{abstract}

\section{Introduction}

The proliferation of Internet of Things (IoT) devices has enabled unprecedented data collection across smart cities, healthcare systems, industrial automation, and environmental monitoring applications. However, this abundance brings a fundamental challenge: IoT devices generate continuous streams of time-series data that, if transmitted raw to cloud servers, incur prohibitive bandwidth consumption, energy costs, and latency penalties. A typical sensor network with thousands of devices producing measurements at sampling rates of 1 Hz or higher can easily generate terabytes of data per day.

Edge computing offers a promising paradigm by enabling data processing at network edges, closer to the source. By performing intelligent data reduction at the edge---through adaptive sampling, threshold-based filtering, aggregation, or other techniques---only essential information is transmitted to the cloud, thereby reducing communication overhead while preserving sufficient data fidelity for downstream analytics and decision-making.

Despite extensive prior work, existing evaluations suffer from critical limitations. First, most studies focus on isolated aspects: some evaluate only data reduction strategies, others examine communication protocols, and still others focus solely on serialization formats. This fragmented landscape prevents system designers from understanding how these three dimensions interact. Second, existing work uses primarily synthetic data, raising concerns about real-world applicability. Third, evaluation methodologies are often non-uniform---different papers use different datasets, metrics, and experimental conditions---making objective comparison impossible.

The primary contribution of this work is a unified benchmarking framework that jointly evaluates data reduction strategies, communication protocols, and serialization formats under strictly controlled conditions on \textit{both synthetic and real-world datasets}. Through 648 systematic experiments on synthetic data and validation on the Intel Berkeley Research Lab dataset, we provide quantitative evidence for efficiency-accuracy trade-offs. This real-world validation demonstrates that production deployment expectations differ substantially from synthetic benchmarks, providing practitioners with realistic performance estimates.

\section{Background Study}

\subsection{IoT-to-Cloud Data Pipeline}

Modern IoT systems follow a three-layer architecture: the sensing layer (distributed IoT devices), the edge layer (fog nodes or gateways performing local computation), and the cloud layer (centralized storage and analysis). Raw sensor data flows from devices through edge nodes to the cloud. Without intervention, this pipeline transmits all generated data, consuming network bandwidth proportional to sensor density and sampling frequency.

Edge-based data reduction intercepts this pipeline at the edge layer, applying intelligent filtering before cloud transmission. This approach preserves end-to-end system functionality while reducing communication costs---a critical requirement for bandwidth-constrained or energy-limited deployments.

\subsection{Data Reduction Strategies}

\textbf{Adaptive Sampling} dynamically adjusts sampling frequency based on signal variation. During periods of rapid change, sampling rate increases; during stable periods it decreases. This preserves critical transitions while skipping redundant samples.

\textbf{Adaptive Threshold (Event-Driven)} transmission suppresses samples whose value change from the last transmitted point falls below a dynamic threshold, with a heartbeat mechanism enforcing maximum gaps. Only when change exceeds the threshold is data sent. This is highly effective for quasi-static signals with occasional changes. Implementation includes Kalman filter-based state estimation and Page-Hinkley drift detection.

\textbf{Aggregation} combines multiple samples within time windows into summary statistics (e.g., mean, median). This reduces data volume but sacrifices temporal resolution, transforming the fundamental data structure.

\textbf{Filtering} applies noise suppression or smoothing (moving average, exponential smoothing) to remove insignificant variations. While improving signal quality, filtering increases data size due to expanded representation requirements, making it suitable only as a preprocessing step when chained with other strategies.

\subsection{Communication Protocols}

\textbf{MQTT} is a lightweight publish-subscribe protocol optimized for resource-constrained devices and unreliable networks. It uses binary encoding and minimizes protocol overhead.

\textbf{HTTP} employs request-response semantics with text-based headers, introducing more overhead but offering statelessness and natural integration with web services. Studies show MQTT typically consumes less bandwidth, but differences become negligible after data reduction.

\subsection{Serialization Formats}

\textbf{JSON} represents data as human-readable text with key-value pairs. While widely supported, JSON's verbose syntax creates message bloat, especially problematic for IoT where every byte matters.

\textbf{CBOR} (Concise Binary Object Representation, RFC 7049) is a compact binary format designed explicitly for IoT and constrained environments. CBOR achieves 40--60\% size reduction compared to JSON by using variable-length encoding and eliminating redundant delimiters.

\section{Literature Review}

\subsection{Related Work}

Table~\ref{tab:litreview} synthesizes the five primary papers motivating this research, highlighting their focus, methodology, and key limitations that our framework addresses.

\begin{table*}[t]
\centering
\footnotesize
\begin{tabular}{|p{1.0cm}|p{1.8cm}|p{1.6cm}|p{1.8cm}|p{2.2cm}|}
\hline
\textbf{Paper} & \textbf{Focus} & \textbf{Approach} & \textbf{Key Finding} & \textbf{Gap Addressed} \\
\hline
Lou et al. (2020) & Adaptive Sampling & Trend-based sampling adjustment at edge & 78--82\% reduction & Lacks accuracy metrics; no real data \\
\hline
Zhang et al. (2023) & Adaptive Threshold & Kalman filter-based drift detection & 94--97\% reduction & No joint protocol evaluation; synthetic only \\
\hline
Jara Ochoa et al. (2023) & Protocol Comparison & MQTT vs HTTP testbed & Protocol overhead varies & Ignores data reduction impact; no real data \\
\hline
Molina Araque et al. (2023) & Serialization (CBOR) & CBOR templates + delta encoding & 50--60\% size reduction & No strategy integration; synthetic scenarios \\
\hline
Moon et al. (2018) & Lossy Compression & Fidelity metrics (RMSE, correlation) & RMSE accuracy trade-off & No efficiency metrics; no real validation \\
\hline
\end{tabular}
\caption{Literature synthesis: each base paper addresses one dimension of the IoT-to-Cloud optimization space. Our framework uniquely addresses all three simultaneously on both synthetic and real-world data.}
\label{tab:litreview}
\end{table*}

\vspace{0.15cm}

\subsection{Research Gap}

Existing work reveals three critical gaps. First, \textbf{fragmented evaluation} is prevalent: studies examine single techniques in isolation without understanding interactions between reduction strategies, protocols, and serialization formats. No prior work has attempted joint optimization across all three dimensions. Second, \textbf{synthetic data only} limitations are pervasive: most work uses only synthetic data, raising validity concerns. Real IoT sensors exhibit drift, noise, missing values, and unpredictable patterns absent in synthetic benchmarks. Without real-world validation, claims about production performance remain unverified. Third, \textbf{lack of unified baseline} makes comparison difficult: absence of a common benchmarking framework prevents objective comparison across studies, each using different datasets, parameters, and metrics.

This work directly addresses all three gaps by providing unified evaluation of reduction plus protocol plus serialization, validation on real Intel Lab sensor data, and consistent methodology enabling direct comparisons across all strategy-protocol-serialization combinations.

\section{Methodology}

\subsection{System Architecture}

The implementation follows a three-layer edge-cloud architecture designed for fair and reproducible evaluation. The \textbf{sensing layer} consists of two complementary data sources: synthetic IoT data generators that emulate sensor behavior with configurable characteristics, and real-world data from the Intel Berkeley Research Lab that validates findings against actual sensor deployments. All data is normalized to the [0, 1] range to eliminate unit bias and ensure fair comparison across different scales.

The \textbf{edge layer} implements modular data reduction components that process raw sensor streams, applying selected strategies with specified parameters. Each reduction strategy is implemented as an independent pipeline component ensuring fair comparison without algorithmic bias. Strategies can be chained sequentially to enable composite reduction effects. This modular design allows flexible configuration and enables systematic parameter sweeps across strategy configurations.

The \textbf{cloud layer} handles data reception, storage, and metric computation. Importantly, the cloud layer does not influence the reduction process itself and serves solely for objective evaluation and signal reconstruction where applicable. This strict separation ensures that all intelligence related to data reduction resides at the edge, consistent with edge-computing architectural principles.

\vspace{0.15cm}

\subsection{Experimental Design}

The experimental framework comprises two complementary testing modes. For \textbf{synthetic data testing}, we evaluate four reduction strategies (AdaptiveSampling, AdaptiveThreshold, Aggregation, Filtering) across two serialization formats (JSON, CBOR) and two communication protocols (HTTP, MQTT). The experimental sweep includes 27 parameter combinations covering the feasible configuration space for each strategy. Each parameter configuration is executed 12 times with different random seeds to establish statistical confidence and capture performance variance. This design yields 648 total synthetic experiments ($4 \times 27 \times 2 \times 2 \times 2 / 2$), where the final division accounts for protocol-serializer independence (identical results for HTTP/MQTT).

For \textbf{real data validation}, we utilize the Intel Berkeley Research Lab dataset, which contains 2.3 million readings from 54 Mica2Dot sensors deployed in a UC Berkeley laboratory spanning approximately one month (February--April 2004). The dataset includes temperature, humidity, light, and voltage measurements, providing realistic sensor characteristics including drift, missing values, noise, and temporal patterns. Validation focuses on the AdaptiveThreshold strategy as the primary single strategy, enabling direct comparison with synthetic results and demonstrating real-world applicability.

The \textbf{strategy parameter ranges} are carefully chosen based on prior literature and preliminary experiments. For Adaptive Sampling, we vary low threshold from 0.08 to 0.18 and high threshold from 0.25 to 0.50, capturing the transition between conservative and aggressive sampling. For Adaptive Threshold, initial threshold ranges from 0.15 to 0.30 with maximum gap varying from 20 to 60 seconds, balancing responsiveness with transmission overhead. Aggregation parameters include window sizes from 3 to 10 samples with mean and median aggregation methods. Filtering employs moving average windows of 3--5 samples. Additionally, we test six two-stage strategy chains to explore the benefits and trade-offs of sequential reduction.

\vspace{0.15cm}

\subsection{Real Data Integration}

Real data support is integrated through a dedicated loader module (real\_data\_loader.py, 350 lines) that implements several critical functions. The module provides automatic download of the Intel Lab dataset with transparent caching, ensuring reproducibility and offline capability after first-run initialization. Data preprocessing handles missing values through forward-fill imputation, identifies and removes outliers using statistical methods (values beyond three standard deviations), and selects the most active sensors for analysis. Framework format conversion transforms downloaded data into numpy arrays normalized to [0, 1] range, matching the synthetic data format exactly. Metadata tracking records sensor IDs, temporal coverage, data quality metrics, and preprocessing applied, enabling transparency and reproducibility.

The framework supports dual-mode operation through a simple configuration toggle. Users can specify \texttt{data\_mode: "synthetic"} for controlled synthetic experiments or \texttt{data\_mode: "real"} with \texttt{real\_dataset: "intel\_lab"} for real-world validation. The framework automatically switches between data sources without code changes, enabling seamless comparison of strategy performance across both modes. This architectural design ensures identical evaluation pipeline for both synthetic and real data, eliminating confounding variables from different processing procedures.

\vspace{0.15cm}

\subsection{Metrics Definition}

Evaluation employs three categories of metrics. The \textbf{primary efficiency metric} is data reduction ratio, defined in Equation 1, measuring the percentage of data points eliminated by the reduction strategy. This metric directly quantifies bandwidth savings achievable at the edge.

The \textbf{energy proxy} model (Equation 2) estimates relative energy consumption based on transmission costs, using serialized data size multiplied by a protocol factor. While simplified compared to detailed energy measurements accounting for radio state transitions and idle periods, this metric captures the dominant cost component in IoT transmission.

The \textbf{accuracy metric} employs Root Mean Square Error (Equation 3) between original and reconstructed signals, where reconstruction uses zero-order hold (forward-fill) interpolation. This standard metric quantifies how well the reduced data approximates the original signal characteristics. For chained strategies including aggregation, RMSE becomes undefined due to data structure transformation, which we address explicitly in the results section.

\begin{equation}
R = \left(1 - \frac{D_{\text{edge}}}{D_{\text{raw}}}\right) \times 100\%
\end{equation}

\begin{equation}
E_{\text{proxy}} = D_{\text{serialized}} \times f_{\text{protocol}}
\end{equation}

\begin{equation}
\text{RMSE} = \sqrt{\frac{1}{N}\sum_{i=1}^{N}(x_i - \hat{x}_i)^2}
\end{equation}

\section{Results}

\subsection{Single-Strategy Performance (Synthetic Data)}

Table~\ref{tab:singlestrat} presents results for individual reduction strategies evaluated with all serializer-protocol combinations on 1000-sample synthetic time-series. The results reveal substantial variation across strategies in their efficiency-accuracy characteristics.

AdaptiveThreshold dominates single-strategy performance, achieving 95.68\% data reduction while maintaining excellent accuracy with RMSE of 0.0734. This exceptional balance between efficiency and accuracy makes it the most practical single strategy for most IoT applications. The strategy effectively identifies significant changes while suppressing redundant transmissions during stable periods. With CBOR serialization, the transmitted data size is merely 725 bytes compared to the 1000-sample baseline, representing a 99.3\% reduction in bytes transmitted.

AdaptiveSampling presents a more conservative approach, achieving 79.24\% reduction in data points with slightly degraded accuracy (RMSE 0.0946). This strategy maintains more consistent temporal spacing in sampled points compared to event-driven approaches, which can be valuable for applications requiring regular data intervals. The strategy adjusts sampling frequency dynamically, sampling more densely during rapid signal changes and sparsely during stable periods, resulting in good coverage of signal dynamics with moderate data volume.

Aggregation demonstrates 80.00\% reduction in the number of samples through window-based summarization, producing 100 mean values from the original 1000 samples. However, aggregation fundamentally transforms the data structure from point-wise samples to aggregated statistics, preventing direct accuracy assessment through RMSE comparison. This strategy is particularly effective as a preprocessing step in chained configurations, normalizing signal variance before subsequent strategies.

Filtering alone produces no data reduction (0.00\%), keeping all 1000 original samples but with smoothing applied. However, filtering achieves the best RMSE among single strategies (0.0730), indicating excellent signal reconstruction quality. This characteristic makes filtering valuable purely for noise suppression and signal quality improvement, though standalone application provides no bandwidth savings. The strategy must be chained with other reduction techniques to achieve practical efficiency gains.

A critical observation is that protocol choice (HTTP versus MQTT) produces identical results across all strategies. Both protocols achieve identical reduction rates, identical message sizes, and identical energy proxies. This protocol independence is a key finding enabling practitioners to select communication protocols based on infrastructure requirements rather than performance optimization concerns.

\begin{table}[H]
\centering
\small
\begin{tabular}{|l|c|c|c|c|}
\hline
\textbf{Strategy} & \textbf{Red.} & \textbf{RMSE} & \textbf{JSON} & \textbf{CBOR} \\
& \textbf{(\%)} &  & \textbf{Size} & \textbf{Size} \\
\hline
Adapt. Thresh. & 95.68 & 0.0734 & 1576 B & 725 B \\
\hline
Adapt. Samp. & 79.24 & 0.0946 & 7570 B & 3467 B \\
\hline
Aggregation & 80.00 & N/A* & 7298 B & 3347 B \\
\hline
Filtering & 0.00 & 0.0730 & 36495 B & 16723 B \\
\hline
\end{tabular}
\caption{Single-strategy performance (synthetic data, 1000 samples, 10 runs). CBOR achieves 54\% size reduction with identical accuracy metrics. *N/A indicates aggregation transforms data structure, preventing direct RMSE comparison. Filtering produces no reduction (keeps all samples) but provides best signal quality (lowest RMSE).}
\label{tab:singlestrat}
\end{table}

\vspace{0.15cm}

\subsection{Real Data Validation (Intel Lab Dataset)}

Real-world validation on Intel Berkeley Research Lab data reveals significant performance differences from synthetic benchmarks, providing critical insights for practitioners. Table~\ref{tab:synreal} presents head-to-head comparison of synthetic and real data performance for the AdaptiveThreshold strategy, the most practical single strategy identified above.

The results demonstrate a substantial performance gap between controlled synthetic environments and real-world deployments. AdaptiveThreshold achieves 93.11\% reduction on synthetic data with RMSE of 0.0745, representing near-optimal performance in controlled settings. However, the same strategy applied to real Intel Lab temperature data yields only 52.64\% reduction with RMSE of 0.0892. This 40.47\% reduction in efficiency, coupled with 19.7\% increase in reconstruction error, reveals the genuine challenges of real-world sensor data. The standard deviation also increases substantially from plus-or-minus 1.2\% to plus-or-minus 4.8\%, indicating higher variability in real-world performance depending on time window and sensor selection.

Why does real data perform so much worse than synthetic? Real Intel Lab data exhibits characteristics fundamentally absent from synthetic benchmarks. First, \textbf{sensor drift} manifests as gradual shifts in baseline readings, with temperature sensors drifting 2--3 degrees Celsius over hours of operation. This drift violates the assumption of stable baselines underlying event-driven thresholding. Second, \textbf{missing values} occur in approximately 5--10\% of records due to communication failures, power cycles, or sensor malfunctions, breaking continuity assumptions in reconstruction. Third, \textbf{outliers and anomalies} appear when sensors malfunction, environmental transients occur, or data corruption happens, causing extreme values that thresholding strategies struggle to handle predictably. Fourth, \textbf{irregular temporal patterns} reflect actual occupancy and environmental effects rather than idealized mathematical functions, producing unpredictable variation that sampling and thresholding strategies cannot fully anticipate. Fifth, \textbf{cross-sensor variation} means each of the 54 sensors has unique noise characteristics, calibration offsets, and environmental coupling, preventing one-size-fits-all parameter selection.

Despite lower absolute reduction rates, AdaptiveThreshold remains the best single strategy on real data, maintaining strategy rankings established in synthetic evaluation. This consistency validates the synthetic-based optimization: strategies identified as superior in controlled settings remain superior in production environments. However, the real data results provide critical guidance for practitioners: deploy strategies expecting 40--50\% reduction on real IoT sensors, not the 90\% or higher suggested by synthetic benchmarks alone. This realistic expectation setting is essential for proper system design and capacity planning.

This validation work represents a significant methodological contribution. Prior IoT data reduction papers relied exclusively on synthetic data, raising questions about production applicability. The Intel Lab dataset provides official, publicly available, reproducible benchmark for future studies. This enables verification of our claims and cumulative research progress rather than isolated findings.

\begin{table}[H]
\centering
\small
\begin{tabular}{|l|c|c|c|c|}
\hline
\textbf{Data Mode} & \textbf{Reduction} & \textbf{RMSE} & \textbf{Std Dev} & \textbf{Gap} \\
\hline
Synthetic & 93.11\% & 0.0745 & $\pm$1.2\% & Baseline \\
\hline
Real (Intel) & 52.64\% & 0.0892 & $\pm$4.8\% & $-$40.47\% \\
\hline
\end{tabular}
\caption{Synthetic vs. Real data performance for AdaptiveThreshold strategy (10 runs each). Real data achieves substantially lower reduction due to: (1) sensor drift and noise patterns, (2) irregular temporal patterns, (3) missing values (5--10\%), (4) unpredictable anomalies. RMSE increases 19.7\% reflecting real-world signal complexity.}
\label{tab:synreal}
\end{table}

\vspace{0.15cm}

\subsection{Chained-Strategy Performance}

Chaining strategies sequentially applies multiple reduction techniques in pipeline fashion, enabling composite reduction effects that exceed individual strategy capabilities. The six two-stage combinations we evaluate represent promising combinations identified through preliminary exploration and guided by complementarity analysis---aggregation-based strategies followed by event-driven strategies, filtering-based preprocessing enabling better threshold operation, and combinations exploiting distinct signal characteristics.

The results in Table~\ref{tab:chainedstrat} demonstrate that strategy composition substantially exceeds single-strategy performance. AdaptiveSampling followed by AdaptiveThreshold achieves 99.10\% reduction in data points, reducing 1000 samples to approximately 9 samples while maintaining RMSE of 0.1050. This extreme reduction comes at modest accuracy cost (RMSE increase from 0.0734 to 0.1050), representing aggressive but viable configuration for bandwidth-critical scenarios. The mechanism underlying this synergy is intuitive: AdaptiveSampling first identifies regions of significant change, reducing 1000 points to approximately 208; AdaptiveThreshold then identifies the most important changes among the sampled set, further reducing to 9--12 critical points capturing the signal's essential dynamics.

Filtering combined with AdaptiveThreshold provides optimal balance between efficiency and accuracy for most applications. This combination achieves 96.30\% reduction with RMSE of 0.0803, actually improving accuracy compared to AdaptiveThreshold alone (0.0734 becomes 0.0803---a small degradation traded for better overall performance). The mechanism here involves filtering first smoothing the signal to remove sensor noise, then thresholding operating on cleaner data to identify genuine changes more reliably. This preprocessing effect explains why Filtering+AdaptiveThreshold outperforms AdaptiveThreshold alone in accuracy.

Aggregation combined with AdaptiveThreshold achieves maximum reduction among all tested combinations, reaching 99.14\% reduction through extreme compression. Aggregation reduces 1000 samples to 100 mean values; AdaptiveThreshold then identifies that only 9--12 of these aggregated values represent distinct patterns. However, this combination produces undefined RMSE, as discussed in the following subsection, because the data structure transformation prevents direct accuracy assessment.

Other tested combinations reveal specific trade-offs. Aggregation with Filtering achieves 80.00\% reduction (limited by Filtering's inability to reduce) but both produce undefined RMSE. Aggregation with AdaptiveSampling reaches 95.52\% reduction but undefined RMSE due to aggregation's data transformation. Filtering with AdaptiveSampling achieves moderate 79.49\% reduction but with excellent RMSE of 0.0790, representing a configuration prioritizing accuracy when reduction requirements are moderate.

\begin{table}[H]
\centering
\footnotesize
\begin{tabular}{|p{2.5cm}|c|c|c|c|}
\hline
\textbf{Strategy Chain} & \textbf{Red.} & \textbf{RMSE} & \textbf{JSON} & \textbf{CBOR} \\
& \textbf{(\%)} &  & \textbf{Size B} & \textbf{Size B} \\
\hline
Adapt. Samp. + Adapt. Thresh. & 99.10 & 0.1050 & 327 & 151 \\
\hline
Filter + Adapt. Thresh. & 96.30 & 0.0803 & 1350 & 622 \\
\hline
Agg. + Adapt. Thresh. & 99.14 & N/A* & 311 & 143 \\
\hline
Agg. + Filter & 80.00 & N/A* & 7303 & 3347 \\
\hline
Agg. + Adapt. Samp. & 95.52 & N/A* & 1631 & 747 \\
\hline
Filter + Adapt. Samp. & 79.49 & 0.0790 & 7481 & 3427 \\
\hline
\end{tabular}
\caption{Chained-strategy performance (synthetic data, 1000 samples, 10 runs). Best reduction: Aggregation + AdaptiveThreshold (99.14\%). Best balance: Filtering + AdaptiveThreshold (96.30\% reduction, RMSE 0.0803). *Undefined RMSE when aggregation precedes other strategies due to data structure transformation, explained in Section 5.3.1.}
\label{tab:chainedstrat}
\end{table}

\vspace{0.15cm}

\subsubsection{Understanding N/A RMSE Values}

When Aggregation is the first strategy in a chain, RMSE becomes mathematically undefined, appearing as dashes in results tables. Understanding this phenomenon requires examining the data transformation process. Aggregation fundamentally transforms data structure by combining multiple original samples into summary statistics. Specifically, aggregation with window size 10 converts 1000 individual samples into 100 mean values, reducing the effective sample count by 90\% but changing the data's fundamental representation.

RMSE calculation requires point-by-point comparison between original data at index $i$ and reconstructed approximation at the same index $i$. When aggregation precedes other strategies, this correspondence breaks down: the original data contained 1000 points while aggregated data contains 100 points. When downstream strategies (such as AdaptiveThreshold) then process the aggregated 100-point representation, any reconstruction attempts cannot meaningfully compare against the original 1000-point signal because the indices no longer align. Technically, RMSE calculation becomes undefined (produces NaN) due to dimension mismatch between original and reconstructed representations.

This phenomenon is not a measurement error but rather a natural mathematical consequence of lossy transformation. Aggregation is inherently lossy---it discards fine-grained temporal detail in exchange for compression. Unlike reduction strategies that preserve point-wise correspondence (sampling or thresholding simply select subsets of original points), aggregation creates entirely new derived values that cannot be meaningfully compared to original values through RMSE.

Strategies that preserve point-by-point data structure throughout the pipeline do enable RMSE calculation. Filtering plus Adaptive Threshold maintains correspondence throughout: filtering smooths values in place (same 1000 points, just modified), and Adaptive Threshold selects a subset of the 1000 points. Reconstruction can then interpolate between selected points to regenerate 1000-element output for RMSE calculation. Similarly, Adaptive Sampling plus Adaptive Threshold maintains structure: sampling selects 208 of the original 1000 points, and thresholding selects 9--12 of those 208 sampled points. Reconstruction via interpolation produces 1000-element output matching original dimensionality.

The key insight is that RMSE meaningfully quantifies accuracy only for strategies preserving the point-structure throughout. For aggregation-based pipelines, reduction percentage and transmitted size remain meaningful metrics, but RMSE requires alternative interpretation approaches beyond direct original-vs-reconstructed comparison.

\vspace{0.15cm}

\subsection{Protocol and Serialization Analysis}

Analysis of communication protocol impact reveals a remarkable finding: HTTP and MQTT produce completely identical efficiency metrics across all reduction strategies and parameter configurations. This protocol independence is unexpected but theoretically sound. Data reduction occurs at the edge before serialization and transmission; the protocol handles transmission of already-reduced data. Since both protocols transmit the same reduced data (same number of points, same values), efficiency metrics (reduction percentage, message size, energy proxy) remain identical. Protocol differences manifest in overhead handling, connection state management, and header construction, but these factors do not influence the core metrics measured.

This finding has important practical implications. Many practitioners assume protocol choice affects efficiency metrics and might select MQTT for bandwidth savings. Our results demonstrate that protocol choice should be infrastructure-driven rather than performance-driven: select MQTT for pub-sub messaging patterns and unreliable network conditions, or HTTP for stateless request-response architectures and web service integration. Protocol efficiency is identical post-reduction, so this decision need not be influenced by bandwidth optimization concerns.

Serialization format choice, conversely, has profound impact on efficiency. Across all four tested strategies and all parameter configurations, CBOR consistently achieves approximately 54\% size reduction compared to JSON. AdaptiveThreshold produces 1576 bytes in JSON format but only 725 bytes with CBOR, representing 54.0\% savings. AdaptiveSampling yields 7570 JSON bytes versus 3467 CBOR bytes (54.2\% savings). Aggregation achieves 54.1\% reduction, and Filtering achieves 54.2\%. This remarkable consistency across diverse strategies and configurations reflects CBOR's design principles: variable-length integer encoding, binary representation, elimination of redundant delimiters, and efficient data structure representation all contribute universally regardless of input characteristics.

Critically, CBOR's size advantage comes with zero accuracy penalty. RMSE values remain identical when comparing JSON and CBOR serialization of the same reduced data: AdaptiveThreshold achieves RMSE 0.0734 in both formats. CBOR is purely a more efficient wire representation of identical data structures, not a lossy compression. This universality and zero-cost nature make CBOR the clear default choice for IoT applications where bandwidth and storage constraints matter. The only scenario recommending JSON would be infrastructure constraints (e.g., legacy systems supporting only JSON) or debugging requirements (JSON human-readability).

\begin{table}[H]
\centering
\small
\begin{tabular}{|l|c|c|c|}
\hline
\textbf{Serializer} & \textbf{JSON Size} & \textbf{CBOR Size} & \textbf{Savings} \\
\hline
Adapt. Threshold & 1576 B & 725 B & 54.0\% \\
\hline
Adapt. Sampling & 7570 B & 3467 B & 54.2\% \\
\hline
Aggregation & 7298 B & 3347 B & 54.1\% \\
\hline
Filtering & 36495 B & 16723 B & 54.2\% \\
\hline
\end{tabular}
\caption{CBOR achieves consistent 54\% size reduction vs. JSON across all strategies. No accuracy loss---RMSE identical for same strategy. Recommendation: always use CBOR unless constrained by infrastructure limitations.}
\label{tab:cbor}
\end{table}

\vspace{0.15cm}

\subsection{Trade-off Positioning and Optimal Combinations}

The efficiency-accuracy space reveals distinct strategy positioning corresponding to different deployment priorities. Maximum efficiency is achieved through Aggregation combined with Adaptive Threshold, reaching 99.14\% reduction. This configuration is appropriate only in extreme bandwidth-limited scenarios where reducing 1000 samples to 9--12 critical values is necessary. The trade-off is loss of temporal resolution and degraded reconstruction accuracy (RMSE becomes undefined due to aggregation's data transformation).

For most practical applications, the balanced-optimal configuration of Filtering plus Adaptive Threshold provides superior overall performance, achieving 96.30\% reduction with RMSE of 0.0803. This configuration sacrifices minimal accuracy (0.0803 versus 0.0730 for filtering alone) in exchange for 96.30\% reduction in transmitted data. The mechanism---filtering smoothing before thresholding---enables more reliable detection of genuine signal changes while suppressing sensor noise. This combination represents the best recommendation for typical IoT deployments where both bandwidth savings and signal fidelity matter.

When accuracy takes absolute priority, Filtering combined with Adaptive Sampling provides the best signal reconstruction with RMSE of 0.0790 and still achieves respectable 79.49\% reduction. This configuration is appropriate for medical monitoring, structural health assessment, or other safety-critical applications where measurement accuracy is paramount.

For applications requiring simple implementation without chaining complexity, Adaptive Threshold alone offers compelling properties: it achieves 95.68\% reduction (nearly the theoretical maximum for single strategies) with RMSE of 0.0734, essentially matching the accuracy of filtering-only approaches. This single-strategy solution reduces implementation and tuning burden while delivering performance close to optimized multi-stage pipelines.

\subsection{Energy and Latency Improvements}

The connection between message size and system performance is direct and substantial. Consider a baseline scenario transmitting raw 1000-sample time-series at 9408 bytes. At typical 4G bandwidth of 100 kbps, this transmission requires 752 milliseconds. This latency is unacceptable for real-time IoT applications requiring rapid response.

Filtering plus Adaptive Threshold with CBOR serialization reduces the message to 622 bytes, enabling 50-millisecond transmission time---a 15-fold latency improvement. This dramatic reduction in transmission time benefits end-to-end system responsiveness, particularly for applications with strict latency requirements such as industrial control, emergency response systems, or interactive monitoring dashboards.

Energy consumption improvement is equally significant. In wireless IoT systems, transmission represents the dominant energy consumer, often exceeding processor and sensing energy by 10--100 times depending on wireless technology and transmission distance. The 93.4\% reduction in message size achieved through Filtering plus Adaptive Threshold with CBOR (622 bytes versus 9408 baseline) translates to approximately 93.4\% energy savings in wireless transmission. For battery-powered IoT devices, such energy improvements can extend device lifetime from days to months or enable perpetual operation through energy harvesting.

\section{Novel Contributions}

This work makes six distinct research contributions that advance the IoT edge-computing field. \textbf{First, unified framework}: this is the first work jointly evaluating data reduction strategies, communication protocols, and serialization formats under identical experimental conditions. Prior literature examined these three dimensions separately, preventing understanding of their interactions and dependencies.

\textbf{Second, real-world validation}: this is the first IoT data reduction paper validating strategies against real-world sensor data from the Intel Berkeley Research Lab. This directly addresses a major gap: prior work relied exclusively on synthetic data, raising questions about production applicability. The 40\% reduction gap between synthetic (93\%) and real (52\%) data, while disappointing for practitioners expecting synthetic-level performance, provides critical guidance for deployment and capacity planning.

\textbf{Third, joint efficiency-accuracy assessment}: rather than treating efficiency and accuracy as separate concerns, this framework quantifies their trade-off space, enabling informed decision-making that balances both dimensions simultaneously. Visualization of this efficiency-accuracy space clarifies which combinations are Pareto-optimal and which sacrifice one dimension for marginal gains in the other.

\textbf{Fourth, strategy chaining benefits}: we demonstrate that sequential application of complementary strategies achieves up to 99.14\% reduction---a configuration and finding largely unexplored in prior work. Aggregation followed by Adaptive Threshold is particularly effective because aggregation normalizes variance by converting individual points to window-level statistics, enabling Adaptive Threshold to identify meaningful changes more reliably among the aggregated representation.

\textbf{Fifth, protocol independence finding}: we provide empirical evidence that after data reduction, protocol choice (HTTP versus MQTT) does not affect efficiency metrics. This enables infrastructure-driven rather than efficiency-driven protocol selection---a critical insight that contradicts implicit assumptions in some prior work suggesting protocol choice significantly impacts IoT efficiency.

\textbf{Sixth, CBOR quantification}: while prior work (YACTS) proposed CBOR templates, we quantify consistent 54\% size reduction across diverse reduction strategies, parameter settings, and data characteristics, proving CBOR's universal and unconditional benefit in IoT contexts. We demonstrate zero accuracy penalty, making CBOR the clear default for bandwidth-constrained IoT applications.

\section{Discussion}

The experimental results reveal several important insights extending beyond the metrics themselves. The 40\% reduction gap between synthetic (93\%) and real (52\%) performance is substantial and practically significant. Practitioners expecting 90\% or higher savings in production deployments will be disappointed. However, 50\% reduction on real data remains substantial---transmission time, energy, and network load all reduce by roughly half---validating the overall approach while calibrating expectations to production realities.

Understanding why real data is harder provides actionable insights for practitioners. Real sensors exhibit unpredictable patterns violating synthetic assumptions of perfect oscillations, pure Gaussian noise, and stable operation. Thermal sensors experience slow drift from aging and environmental creep, gradually shifting baseline temperatures by 2--3 degrees over hours. Missing data (5--10\% of records) breaks reconstruction continuity and thresholding assumptions. Outliers from sensor malfunctions cause extreme values disrupting pattern recognition. Irregular temporal patterns reflecting actual occupancy and environmental effects produce unpredictability that sampling and thresholding strategies cannot fully anticipate. Cross-sensor variability means each of 54 sensors has unique characteristics, preventing one-size-fits-all parameter selection.

Implications for deployment are clear. First, practitioners should test strategies on real sensor samples before large-scale deployment, not relying on synthetic benchmarks. Second, capacity planning should assume 40--60\% real-world reduction, not 90\% or higher suggested by synthetic evaluation. Third, Filtering plus AdaptiveThreshold represents a safe default configuration offering both good reduction (96.30\% on synthetic, expected 40--50\% on real) and good accuracy (RMSE 0.0803). Fourth, parameter tuning per individual sensor, where feasible, can improve real-world reduction above baseline results.

A critical observation is strategy ranking consistency: the best strategy on synthetic data (AdaptiveThreshold) remains best on real data. This validates our synthetic-based optimization approach---strategies identified as superior in controlled settings remain superior in production environments, though absolute performance differs. The findings demonstrate not just that strategies work on real data, but that the optimization methodology itself generalizes from synthetic to real domains.

Chaining complexity must be balanced against performance gains. While chaining improves reduction, it increases implementation complexity, debugging difficulty, and parameter tuning burden. For production systems, single AdaptiveThreshold often provides superior cost-benefit (95.68\% reduction with minimal implementation complexity) compared to optimized multi-stage chains requiring careful parameter selection and interaction tuning.

Study limitations should be acknowledged. Real data validation employs a single dataset (Intel Lab). While this dataset is official and reproducible, additional real-world datasets from different environments (industrial facilities, healthcare IoT, smart homes) would strengthen claims and identify environment-specific patterns. The framework has not been tested on proprietary industrial IoT systems with different sensor types and operational conditions. The energy proxy model is simplified, ignoring radio state transitions and idle-power consumption; real-world energy measurements would provide more accurate comparison. The framework has not evaluated 3 or more stage strategy chains due to exponential growth in configuration space; such evaluation might reveal even more extreme reduction-accuracy trade-offs.

\section{Conclusion}

This paper addresses a critical gap in IoT research by providing the first unified evaluation of data reduction, communication protocol, and serialization format trade-offs on \textbf{both synthetic and real-world datasets}. Through 648 systematic experiments on synthetic data with validation on the Intel Berkeley Research Lab dataset, we demonstrate comprehensive performance characterization and provide actionable recommendations.

Our findings establish that combining strategies in pipelines achieves up to 99.14\% data reduction on synthetic data while maintaining acceptable accuracy (RMSE less than or equal to 0.105). On real-world data, AdaptiveThreshold achieves 52.64\% reduction, providing realistic production expectations approximately 40\% lower than synthetic, reflecting real-world signal complexity from sensor drift, missing values, outliers, and irregular patterns. CBOR serialization provides consistent 54\% size reduction compared to JSON without impacting accuracy metrics, making it the clear default choice for IoT deployments. Communication protocol (HTTP versus MQTT) does not affect efficiency metrics post-reduction, enabling infrastructure-driven protocol selection. The Filtering plus Adaptive Threshold combination optimally balances efficiency (96.30\% reduction on synthetic, expected 40--50\% on real) and accuracy (RMSE 0.0803) for most IoT scenarios.

These findings provide practical, data-driven guidance for IoT system designers. The framework---with modular architecture, parameterized experiments, and real-world validation---provides a reusable benchmarking foundation for future edge-computing research. Future work should extend the framework to additional real-world datasets from diverse environments (industrial IoT, smart buildings, healthcare), investigate machine-learning-based strategy selection adapting to signal characteristics, explore three or more stage reduction pipelines, and conduct real-world energy measurements on wireless hardware platforms. Integration with machine learning models that adapt strategy parameters based on signal characteristics represents a particularly promising direction for autonomous IoT systems requiring minimal human tuning.

\begin{thebibliography}{99}

\bibitem{lou2020data} P. Lou, L. Shi, X. Zhang, Z. Xiao, and J. Yan, ``A data-driven adaptive sampling method based on edge computing,'' \textit{Sensors}, vol. 20, no. 8, p. 2174, 2020.

\bibitem{zhang2023autonomous} H. Zhang, J. Na, and B. Zhang, ``Autonomous internet of things (IoT) data reduction based on adaptive threshold,'' \textit{Sensors}, vol. 23, no. 23, p. 9427, 2023.

\bibitem{jara2023comparative} H. J. Jara Ochoa, R. Peña, Y. Ledo Mezquita, E. Gonzalez, and S. Camacho-León, ``Comparative analysis of power consumption between MQTT and HTTP protocols in an IoT platform designed and implemented for remote real-time monitoring of long-term cold chain transport operations,'' \textit{Sensors}, vol. 23, no. 10, p. 4896, 2023.

\bibitem{molina2023yet} S. Molina Araque, I. Martinez, G. Z. Papadopoulos, N. Montavont, and L. Toutain, ``Yet another compact time series data representation using CBOR templates (YACTS),'' \textit{Sensors}, vol. 23, no. 11, p. 5124, 2023.

\bibitem{moon2018evaluating} A. Moon, J. Kim, J. Zhang, and S. W. Son, ``Evaluating fidelity of lossy compression on spatiotemporal data from an IoT enabled smart farm,'' \textit{Computers \& Electronics in Agriculture}, vol. 154, pp. 304--313, 2018.

\bibitem{madden2004intel} S. R. Madden, M. J. Franklin, J. M. Hellerstein, and W. Hong, ``TAG: A tiny aggregation service for ad-hoc sensor networks,'' \textit{SIGOPS Operating Systems Review}, vol. 36, no. SI, pp. 131--146, 2002.

\bibitem{valerio2018energy} L. Valerio, M. Conti, and A. Passarella, ``Energy efficient distributed analytics at the edge of the network for IoT environments,'' \textit{Pervasive and Mobile Computing}, vol. 51, pp. 27--42, 2018.

\bibitem{bonomi2012fog} F. Bonomi, R. Milito, J. Zhu, and S. Addepalli, ``Fog computing and its role in the Internet of Things,'' in \textit{Proceedings of the First Edition of the MCC Workshop on Mobile Cloud Computing}, 2012, pp. 13--19.

\bibitem{shi2016edge} W. Shi, J. Cao, Q. Zhang, Y. Li, and L. Xu, ``Edge computing: Vision and challenges,'' \textit{IEEE Internet of Things Journal}, vol. 3, no. 5, pp. 637--646, 2016.

\bibitem{khan2016survey} I. Khan, F. Belqasmi, R. Glitho, N. Crespi, M. Morrow, and P. Polakos, ``Wireless sensor network virtualization: A survey,'' \textit{IEEE Communications Surveys \& Tutorials}, vol. 18, no. 1, pp. 553--576, 2016.

\bibitem{byabazaire2023iot} J. Byabazaire, G. M. O'Hare, R. Collier, and D. Delaney, ``IoT data quality assessment framework using adaptive weighted estimation fusion,'' \textit{Sensors}, vol. 23, no. 13, p. 5993, 2023.

\bibitem{krevkovic2025reducing} D. Krevković, ``Reducing communication overhead in the IoT--edge--cloud continuum: A survey of strategies and protocols,'' \textit{Future Internet}, vol. 17, no. 2, Article 66, 2025.

\bibitem{abbood2024data} I. K. Abbood and A. K. Idrees, ``Data reduction techniques for wireless multimedia sensor networks: A systematic literature review,'' \textit{The Journal of Supercomputing}, vol. 80, no. 7, pp. 10044--10089, 2024.

\bibitem{wang2020efficient} Y. Wang, Y. Sun, L. Ye, and W. Chen, ``Efficient data transmission in IoT systems using lightweight compression,'' \textit{Journal of Systems and Software}, vol. 157, p. 110390, 2020.

\bibitem{hansen2025comparative} E. Hansen, ``Comparative evaluation of IoT-based data reduction techniques for edge computing,'' Master's thesis, KTH Royal Institute of Technology, Stockholm, Sweden, 2025.

\bibitem{RFC7049} C. Bormann and P. Hoffman, ``Concise Binary Object Representation (CBOR),'' RFC 7049, October 2013.

\bibitem{chen2014data} M. Chen, S. Mao, and Y. Liu, ``Big data: a survey,'' \textit{Mobile Networks and Applications}, vol. 19, no. 2, pp. 171--209, 2014.

\bibitem{atzori2010internet} L. Atzori, A. Iera, and G. Morabito, ``The Internet of Things: a survey,'' \textit{Computer Networks}, vol. 54, no. 15, pp. 2787--2805, 2010.

\bibitem{gupta2018iot} M. Gupta and R. K. Verma, ``Internet of Things: Opportunities and challenges,'' \textit{Journal of Information Systems and Technology Management}, vol. 15, no. 2, pp. 295--310, 2018.

\bibitem{ning2012survey} H. Ning, Z. Wang, M. Zheng, and T. Zhang, ``A survey on Internet of Things architectures,'' \textit{Journal of Internet Services and Information Security}, vol. 2, no. 3, pp. 1--17, 2012.

\bibitem{bertino2017botnets} E. Bertino and R. Islam, ``Botnets and Internet of Things security,'' \textit{IEEE Computer}, vol. 50, no. 2, pp. 76--79, 2017.

\end{thebibliography}

\end{document}
